\section{Fase 2: Especificação da Avaliação --- Objetivos, Questões e Hipóteses}
\label{sec:fase2-gqm}

Em conformidade com o processo de avaliação da ISO/IEC 25040 \cite{iso25040_2011} e com a abordagem \textit{Goal Question Metric} (GQM) \cite{basili1994gqm}, a presente seção operacionaliza as características priorizadas na Fase 1, transformando-as em objetivos formais de medição, questões analíticas orientadoras e suas correspondentes hipóteses de teste.

\subsection{Rastreabilidade entre a Fase 1 e a Fase 2}
\label{subsec:rastreabilidade}

A definição dos elementos da Fase 2 preserva estrita rastreabilidade com as decisões tomadas na Fase 1. Cada objetivo GQM possui seu foco de qualidade derivado diretamente das características de qualidade priorizadas na etapa anterior: Adequação Funcional (ISO/IEC 25010:2023), Eficiência de Desempenho (ISO/IEC 25010:2023) e Cobertura de Contexto (ISO/IEC 25010:2011 e ISO/IEC 25022:2016). 

Da mesma forma, o ponto de vista de cada objetivo ancora-se nas partes interessadas identificadas no mapeamento de \textit{stakeholders}:
\begin{itemize}
    \item O \textbf{Cidadão/Condutor} subsidia o ponto de vista prioritário sobre a adequação funcional e a cobertura de contexto, visto que necessita de dados íntegros e acesso contínuo independentemente do dispositivo ou da infraestrutura de telecomunicações disponível;
    \item A \textbf{Equipe Técnica de TI e Gestão de Serviços Digitais do Detran-DF} orienta o ponto de vista de eficiência de desempenho, cujo interesse recai sobre a otimização de banda e a estabilidade dos tempos de execução da consulta;
    \item O \textbf{Requisitante da Avaliação (Área de Governança Digital do Detran-DF)} fundamenta transversalmente a busca pela paridade e consistência entre os canais de atendimento prestados pelo órgão.
\end{itemize}

A Tabela~\ref{tab:rastreabilidade-fase1-fase2} sintetiza as relações de rastreabilidade entre as fases.

\begin{table}[htbp]
\centering
\caption{Matriz de Rastreabilidade entre a Fase 1 e a Fase 2 (GQM)}
\label{tab:rastreabilidade-fase1-fase2}
\footnotesize
\begin{tabular}{|p{2.8cm}|p{3.2cm}|p{2.8cm}|p{2.8cm}|p{1.8cm}|}
\hline
\textbf{Característica (Fase 1)} & \textbf{Subcaracterísticas Consideradas} & \textbf{Foco de Qualidade (GQM)} & \textbf{Ponto de Vista (\textit{Stakeholder})} & \textbf{Objetivo GQM} \\ \hline
Adequação Funcional (ISO/IEC 25010:2023) & Completude funcional, correção funcional, adequação funcional & Completude, correção e adequação do fluxo de consulta & Cidadão / Condutor e Requisitante & \textbf{G1} \\ \hline
Eficiência de Desempenho (ISO/IEC 25010:2023) & Comportamento temporal, utilização de recursos & Tempo até a exibição completa da lista de infrações e volume de dados transferidos via rede & Equipe Técnica de TI e Gestão de Serviços Digitais do Detran-DF & \textbf{G2} \\ \hline
Cobertura de Contexto (Qualidade em Uso) & Completude de contexto, flexibilidade operacional & Sucesso da tarefa em condições diversas de hardware e conectividade & Cidadão / Condutor e Requisitante & \textbf{G3} \\ \hline
\end{tabular}
\end{table}

\subsection{Objetivos de Medição no Modelo Estruturado GQM}
\label{subsec:objetivos-gqm}

Para cada uma das características selecionadas, formulou-se um objetivo de medição formal seguindo rigorosamente o \textit{template} de cinco dimensões proposto por Basili e Rombach \cite{basili1994gqm}. A Tabela~\ref{tab:objetivos-gqm} especifica os objetivos concebidos para a avaliação comparativa dos canais digitais do Detran-DF.

\begin{table}[htbp]
\centering
\caption{Objetivos de Medição no Padrão Estruturado GQM}
\label{tab:objetivos-gqm}
\footnotesize
\begin{tabular}{|l|p{12.2cm}|}
\hline
\textbf{ID} & \textbf{Especificação Estruturada do Objetivo} \\ \hline
\textbf{G1} & \textbf{Analisar} o serviço de consulta de infrações de trânsito \newline
\textbf{com o propósito de} comparar e avaliar \newline
\textbf{com respeito à} adequação funcional (completude funcional, correção funcional e adequação funcional do fluxo de consulta) \newline
\textbf{do ponto de vista do} cidadão/condutor e do requisitante institucional \newline
\textbf{no contexto da} execução da consulta por meio do Portal de Serviços web e do aplicativo mobile oficial do Detran-DF com veículos reais registrados no DF. \\ \hline
\textbf{G2} & \textbf{Analisar} o serviço de consulta de infrações de trânsito \newline
\textbf{com o propósito de} comparar e avaliar \newline
\textbf{com respeito à} eficiência de desempenho (tempo até a exibição completa da lista de infrações e volume de dados transferidos via rede) \newline
\textbf{do ponto de vista da} Equipe Técnica de TI e Gestão de Serviços Digitais do Detran-DF e do requisitante institucional \newline
\textbf{no contexto de} requisições repetidas de consulta executadas em ambientes controlados sob condições de rede estável e degradada. \\ \hline
\textbf{G3} & \textbf{Analisar} o serviço de consulta de infrações de trânsito \newline
\textbf{com o propósito de} comparar e avaliar \newline
\textbf{com respeito à} cobertura de contexto (capacidade de conclusão da tarefa mandatória de consulta frente à diversidade de ambientes operacionais) \newline
\textbf{do ponto de vista do} cidadão/condutor sob restrições técnicas \newline
\textbf{no contexto de} variação de plataformas de acesso (sistemas operacionais Android/iOS, navegadores web modernos) e perfis de conectividade (banda larga Wi-Fi, dados móveis 4G e conexões com latência simulada). \\ \hline
\end{tabular}
\end{table}

\subsection{Questões de Medição e Hipóteses}
\label{subsec:questoes-hipoteses}

A derivação do nível operacional desdobra cada objetivo em questões analíticas que abordam os atributos essenciais de qualidade do produto, deliberadamente evitando formular métricas diretas em formato interrogativo. Para cada questão formulada, estabeleceu-se uma hipótese verificável por meio das medições definidas para a Fase 3.
\subsubsection{Desdobramento do Objetivo G1 (Adequação Funcional)}

O foco de adequação funcional investiga se ambos os canais entregam ao cidadão as informações necessárias à identificação e compreensão da autuação, conforme lista de referência definida na Fase 1, e se mantêm consistentes entre si, além de permitirem concluir a consulta sem passos desnecessários.

\begin{itemize}
    \item \textbf{Questão Q1.1 (Completude Funcional):} Os canais web e mobile apresentam as informações necessárias à identificação e compreensão da autuação de trânsito, conforme lista de referência definida na Fase 1?
    \begin{quote}
        \textbf{Hipótese H1.1:} O Portal de Serviços web apresenta maior detalhamento e completude dos campos informativos definidos na lista de referência da Fase 1 do que o aplicativo mobile.
    \end{quote}
    
    \item \textbf{Questão Q1.2 (Correção e Consistência Funcional):} As informações referentes ao status, valores monetários, pontuação e prazos das autuações apresentam consistência observável dos campos comuns entre os canais?
    \begin{quote}
        \textbf{Hipótese H1.2:} Não há divergências observáveis entre os canais para uma mesma infração consultada em intervalo temporal próximo.
    \end{quote}

    \item \textbf{Questão Q1.3 (Adequação das Operações):} A sequência operacional exigida para atingir o extrato final de infrações impõe etapas redundantes ou passos transacionais desnecessários em algum dos canais?
    \begin{quote}
        \textbf{Hipótese H1.3:} O aplicativo mobile proporciona um caminho operacional mais enxuto (menor número de cliques e transições de tela) até a exibição das infrações do que o Portal de Serviços web.
    \end{quote}
\end{itemize}

\subsubsection{Desdobramento do Objetivo G2 (Eficiência de Desempenho)}

O objetivo de eficiência de desempenho avalia o tempo necessário para a conclusão da consulta e o volume de dados transferidos, respeitando a premissa de medição em caixa-preta:

\begin{itemize}
    \item \textbf{Questão Q2.1 (Comportamento Temporal):} Qual canal apresenta menor tempo decorrido até a exibição completa da lista de infrações?
    \begin{quote}
        \textbf{Hipótese H2.1:} O aplicativo mobile apresenta tempo médio inferior ao do Portal de Serviços web para exibir completamente a lista de infrações, nas mesmas condições de rede e para o mesmo veículo consultado.
    \end{quote}

    \item \textbf{Questão Q2.2 (Utilização de Recursos de Rede):} Qual canal transfere menor volume de dados para concluir a consulta de infrações?

    \begin{quote}
        \textbf{Hipótese H2.2:} O aplicativo mobile apresenta menor volume médio de dados transferidos via rede na transação de consulta quando comparado ao Portal de Serviços web.
    \end{quote}

\end{itemize}

\subsubsection{Desdobramento do Objetivo G3 (Cobertura de Contexto)}

Embasada no modelo de qualidade em uso (e distinguindo-se terminantemente de aspectos subjetivos de usabilidade, os quais foram expressamente excluídos deste estudo), a cobertura de contexto afere a eficácia de conclusão da tarefa sob múltiplos ecossistemas de software, hardware e infraestrutura física:

\begin{itemize}
    \item \textbf{Questão Q3.1 (Completude de Contexto Tecnológico):} A funcionalidade de consulta de infrações é concluída com sucesso em todos os ambientes operacionais definidos no escopo da avaliação (combinações de sistemas operacionais e navegadores)?
    \begin{quote}
        \textbf{Hipótese H3.1:} O Portal de Serviços web e o aplicativo mobile permitem a conclusão integral da tarefa de consulta em 100\% dos ambientes de teste homologados (desktop e mobile), sem bloqueios impeditivos decorrentes da plataforma.
    \end{quote}

    \item \textbf{Questão Q3.2 (Flexibilidade em Condições de Rede Degradada):} A consulta de infrações continua sendo concluída quando os canais são submetidos a condições de rede degradada, em comparação com as condições de referência Wi-Fi e 4G?
    \begin{quote}
        \textbf{Hipótese H3.2:} A taxa de conclusão da consulta de infrações é menor sob conexões com latência simulada do que sob Wi-Fi e 4G em pelo menos um dos canais.
    \end{quote}
\end{itemize}

A Tabela~\ref{tab:resumo-questoes-hipoteses} consolida o arcabouço conceitual e operacional derivado, servindo de insumo formal e direto para a especificação das métricas quantitativas, níveis de julgamento e do diagrama GQM a serem desenvolvidos na continuidade da Fase 2.

\begin{table}[htbp]
\centering
\caption{Quadro Sinóptico de Objetivos, Questões e Hipóteses de Medição}
\label{tab:resumo-questoes-hipoteses}
\footnotesize
\begin{tabular}{|c|p{3.0cm}|c|p{4.2cm}|p{4.5cm}|}
\hline
\textbf{Objetivo} & \textbf{Foco de Qualidade} & \textbf{ID} & \textbf{Questão de Medição} & \textbf{Hipótese Formulada} \\ \hline
\multirow{3}{*}{\textbf{G1}} & \multirow{3}{=}{Adequação Funcional (ISO/IEC 25010:2023)} & \textbf{Q1.1} & Integralidade dos dados do auto de infração entre os canais & \textbf{H1.1:} Portal web apresenta maior completude de dados detalhados que o app. \\ \cline{3-5}
& & \textbf{Q1.2} & Consistência observável dos campos comuns entre os canais & \textbf{H1.2:} Não há divergências observáveis nos campos comuns entre os canais. \\ \cline{3-5}
& & \textbf{Q1.3} & Número de etapas até a exibição da lista de infrações & \textbf{H1.3:} Aplicativo mobile requer menos passos operacionais que o portal. \\ \hline
\multirow{2}{*}{\textbf{G2}} & \multirow{2}{=}{Eficiência de Desempenho (ISO/IEC 25010:2023)} & \textbf{Q2.1} & Tempo até a exibição completa da lista de infrações & \textbf{H2.1:} Aplicativo mobile possui menor tempo de resposta que o portal web. \\ \cline{3-5}
& & \textbf{Q2.2} & Volume total de dados transferidos via rede na transação & \textbf{H2.2:} Aplicativo consome menor volume de tráfego de rede que o portal. \\ \hline
\multirow{2}{*}{\textbf{G3}} & \multirow{2}{=}{Cobertura de Contexto (Qualidade em Uso)} & \textbf{Q3.1} & Conclusão da tarefa na matriz de plataformas homologadas & \textbf{H3.1:} Ambos os canais concluem a consulta em 100\% dos ambientes-alvo. \\ \cline{3-5}
& & \textbf{Q3.2} & Operabilidade e resiliência sob redes de alta latência e baixa banda & \textbf{H3.2:} Taxa de conclusão é menor sob latência simulada do que sob Wi-Fi e 4G em pelo menos um dos canais. \\ \hline
\end{tabular}
\end{table}
